\chapter{INTRODUÇÃO}

A computação gráfica tem evoluído significativamente nas últimas décadas, impulsionada pela necessidade de gerar imagens cada vez mais realistas em diferentes contextos. Nesse cenário, o uso de algoritmos de \textit{ray tracing} tem se destacado pela sua capacidade de produzir imagens altamente realistas. No entanto, o custo computacional associado a essa técnica ainda é um desafio, especialmente em cenários dinâmicos.

Para viabilizar o \textit{ray tracing} em tempo real, torna-se fundamental o uso de estruturas de aceleração. As \glspl{bvh} permitem reduzir significativamente o número de testes de interseção entre raios e primitivas geométricas. Com o avanço do hardware dedicado, como \gls{gpu} moderna, diferentes estratégias de construção e atualização dessas estruturas têm sido propostas, cada uma apresentando características distintas em termos de desempenho e eficiência.

Diante desse cenário, torna-se relevante investigar como essas diferentes abordagens se comportam em aplicações de tempo real, considerando tanto as características das cenas quanto as limitações do hardware utilizado.

\section{Problema de Pesquisa}

Com o avanço de aplicações em tempo real, como renderização gráfica e simulações interativas, torna-se essencial o uso de estruturas de aceleração eficientes, como as variantes de \gls{bvh}. No entanto, diferentes algoritmos e estratégias de construção e atualização dessas estruturas, como o \textit{rebuild} completo e o \textit{refit} incremental, apresentam comportamentos distintos dependendo do tipo de hardware utilizado e das características da cena.

Diante disso, surge o seguinte problema de pesquisa: qual abordagem de construção e atualização de \gls{bvh} apresenta melhor desempenho em simulações de tempo real, considerando diferentes cenários de hardware dedicado e cenas dinâmicas?

\section{Objetivos}

\subsection{Objetivo Geral}

Analisar e comparar diferentes estruturas de aceleração baseadas em \gls{bvh}, avaliando seu desempenho em simulações de tempo real, quando utilizadas em conjunto com hardware dedicado.

\subsection{Objetivos Específicos}

\begin{itemize}
    \item Analisar as vantagens e desvantagens de diferentes estruturas \gls{bvh} em simulações de tempo real;
    \item Avaliar a relevância das estruturas de aceleração em cenários com hardware dedicado;
    \item Investigar abordagens de construção de \gls{bvh}, incluindo \textit{rebuild} completo e \textit{refit} incremental, em cenas dinâmicas e semi-dinâmicas;
    \item Comparar diferentes estruturas \gls{bvh} considerando aspectos:
    \begin{itemize}
        \item latência por frame;
        \item uso de memória;
        \item ocupação da \gls{gpu} e \gls{cpu};
        \item consumo energético médio;
    \end{itemize}
\end{itemize}

\section{Justificativa}

A utilização de estruturas de aceleração, como \gls{bvh}, é fundamental para viabilizar aplicações de alto desempenho em computação gráfica, especialmente em contextos de tempo real. Com a crescente adoção de hardware dedicado, como GPUs modernas e unidades específicas para \textit{ray tracing}, torna-se necessário compreender como diferentes estratégias de construção e atualização dessas estruturas impactam o desempenho geral do sistema.

Além disso, abordagens como \textit{rebuild} completo e \textit{refit} incremental apresentam características distintas, que podem influenciar diretamente a eficiência da aplicação dependendo do cenário de uso. A análise desses fatores é importante não apenas do ponto de vista teórico, mas também prático, contribuindo para a escolha de técnicas mais adequadas em aplicações reais.

Dessa forma, este trabalho se justifica pela necessidade de avaliar e comparar essas abordagens, considerando múltiplas métricas de desempenho, a fim de fornecer subsídios para o desenvolvimento de sistemas mais eficientes em ambientes de tempo real.
